\chapter{Planning periods}\label{ch:periods}

A planning period is a window of days (usually two weeks) that is scheduled as one problem. It takes the
factory set-up and the open orders; you check its inputs, generate and commit a schedule, run it, and roll
over to the next period.

\section{The list of periods}

\shot{periods}{Planning periods: one open period (with its variant and reschedule) and one closed period.}{periods}

The page groups periods into \ui{Current \& upcoming} and \ui{Closed}. For each period you see its name and
dates, its status with the next step as a link, the number of orders inside and after it, the number of
options (and how many are still solving), the committed option and when it was committed, and the last
update. What-if variants and reschedules are listed under their main period. Click a period name to open
it; click the column headers to sort.

\subsection{Period statuses}\label{sec:statuses}

\begin{center}
\begin{tabularx}{\linewidth}{>{\bfseries}p{3.7cm}Xp{4.3cm}}
\toprule
\rowcolor{white}Status & Meaning & Next step shown \\
\midrule
Draft & Created, inputs not checked yet. & Check orders and opening stock \\
Inputs need fixing & The last save found errors; schedules cannot be generated. & Fix the errors under Prepare \\
Ready to schedule & Inputs are valid. & Generate schedule options \\
Solving & Runs are queued or running. & Wait for the solver \\
Options ready & At least one feasible schedule exists. & Compare and commit a schedule \\
Committed & An option has been committed. & Issue work orders, watch for incidents \\
Closed & The period is finished and read-only. & Period finished \\
\bottomrule
\end{tabularx}
\end{center}

A period whose end date has passed while it is still open is flagged \ui{Overdue to close}: its undelivered
orders are stuck in it until it is closed.

\section{Creating a period}\label{sec:new-period}

Click \ui{Plan a new period} (on this page, on Today, or in the \ui{New} menu). You only pick the dates; plant
set-up, shifts, rest days and open orders come in automatically.

\shotf[0.55\linewidth]{new_period}{Plan a new period. The warnings explain consequences before the period is created.}{new_period}

\begin{description}[style=nextline, leftmargin=1.2em, font=\bfseries\color{brand}]
  \item[Starts] Proposed as the day after the last period (``Right after Weeks 48-49''). A date in the past is
        allowed but flagged: the period would schedule past days.
  \item[Days] Defaults to the factory's period length. The resulting date range is shown below.
  \item[Name (optional)] Left empty, the period is named after its dates.
  \item[Opening stock] \ui{From the previous committed schedule}: finished goods, semi-finished stock and mold
        counters are projected from the closing state of the previous committed period, and you only correct
        what differs. \ui{Start from zero and enter counts}: for the very first period or after a stock-take.
  \item[Take open orders from the order book] Shows how many orders fall in the period and how many in the
        look-ahead window after it.
\end{description}

The dialog warns you about situations that usually indicate a mistake:

\begin{itemize}
  \item orders already due before the start --- they are taken into the period to ship as early as possible;
        consider telling the customers;
  \item the previous period is committed but not closed --- the orders it delivers are still open and would
        be taken again. Prefer opening it, going to Execute and using \ui{Close \& start the next period}
        (Section~\ref{sec:close}), which closes it and creates the next period in one step;
  \item the new period overlaps an open period --- allowed, but both may schedule the same orders.
\end{itemize}

Click \ui{Create period}. The period opens on its Prepare step.

\section{The period page}\label{sec:period-page}

\shot{period_header}{The period page of ``Weeks 48-49'' (here on the Prepare step).}{period_header}

\begin{callouts}
  \item[1] \textbf{Name and status.} Variants and reschedules add their label after the name and a kind badge
           (\ui{What-if variant}, \ui{Reschedule}).
  \item[2] \textbf{Key facts}: dates, length, orders in and after the period, and the factory set-up version
           the period uses (click it to open the Factory page). A reschedule also shows ``fixed before'' its
           cut-off.
  \item[3] \textbf{Versions of the period}: the main period and its reschedules and variants as chips; the
           committed one has a tick. Click a chip to switch (Chapter~\ref{ch:variants}).
  \item[4] \textbf{Next-step banner}: what to do now, with a button that goes there --- for example
           ``2 options are ready. Next: compare them and commit one.'' or, as here, ``Committed \dots\ Next:
           issue work orders, track incidents, and close the period at the end.''
  \item[5] \textbf{Step tabs}: Prepare, Options, Schedule \& compare, Commit, Execute. A green tick marks a done
           step, a number the current or a future step, a red dot a step with errors.
  \item[6] \textbf{Analysis and History}: how the scheduling went, and the change log (Chapter~\ref{ch:analysis}).
  \item[7] \textbf{More} and \textbf{Generate options}/\textbf{Generate more}.
\end{callouts}

\shotf[0.32\linewidth]{more_menu}{The \ui{More} menu of a period.}{more_menu}

The \ui{More} menu offers \ui{Create a what-if variant\dots} (Chapter~\ref{ch:variants}), \ui{Rename\dots}
and \ui{Delete period\dots} (for a variant or reschedule: \ui{Delete this version\dots}). Deleting removes
the inputs, options and decisions of that period or version and of its child variants and reschedules; the
order book is not affected.

\begin{note}
The tab you are on is part of the address (for example \code{?tab=schedule}), so you can bookmark a step or
send a colleague a link to it.
\end{note}
